\documentclass{article}
\usepackage[utf8]{inputenc}
\usepackage[T1]{fontenc}
\usepackage{geometry}
\usepackage{graphicx}
\usepackage{float}
\usepackage{hyperref}
\usepackage[english,indonesian]{babel}

\title{Optimasi Pemrosesan Template LaTeX Menggunakan Pendekatan Berbasis AI}
\author{Academic Support Agent}
\date{\today}

\begin{document}
\maketitle

\begin{abstract}

\end{abstract}

\chapter{PENDAHULUAN}


\chapter{TINJAUAN PUSTAKA}


% Original content (commented out):
% \documentclass[12pt,a4paper]{article}
\title{Plantilla para trabajos UPIITA V1}
\usepackage{UPIITA}
\asignatura{Asignatura}
\carrera{Ing. Mecatrónica}
\titulo{Título}
\grupo{Grupo}
\alumno{Alumno}
\profesor{Profesor}
\begin{document}
\maketitlepage % Portada
\newpage
\tableofcontents %Índice
\newpage
\listoffigures   %Índice de figuras
%\newpage        
%\listoftables   %Índice de tablas

% Cuerpo del Documento
\newpage
\section{Section}
\section{Section} 

% Insertar .tex al documento
\input{ParteDelDocumento}

% Bibliografia o referencias
\newpage
\addcontentsline{toc}{section}{Bibliografía}
\printbibliography[title={Bibliografía}] %Referencias

\end{document}

\end{document}